%% ============================================================
%%  The dm3 Operator: Explicit Toy Model and
%%  Global Dynamical Analysis
%%
%%  Paper 2 of 2 (companion to Paper 1)
%%  Target: SIAM Journal on Applied Dynamical Systems / Chaos / Physica D
%% ============================================================

\documentclass[reqno]{amsart}

\usepackage{amsmath,amssymb,amsthm,amsfonts}
\usepackage{mathtools}
\usepackage{hyperref}
\usepackage{enumitem}
\usepackage{booktabs}
\usepackage{array}

\hypersetup{colorlinks=true,linkcolor=blue,citecolor=blue,urlcolor=blue}

\numberwithin{equation}{section}

\theoremstyle{plain}
\newtheorem{theorem}{Theorem}[section]
\newtheorem{proposition}[theorem]{Proposition}
\newtheorem{corollary}[theorem]{Corollary}
\newtheorem{lemma}[theorem]{Lemma}

\theoremstyle{definition}
\newtheorem{definition}[theorem]{Definition}
\newtheorem{axiom}{Axiom}

\theoremstyle{remark}
\newtheorem{remark}[theorem]{Remark}

%% Macros (matching Paper 1)
\newcommand{\R}{\mathbb{R}}
\newcommand{\T}{\mathbb{T}}
\newcommand{\dm}{\mathfrak{D}}
\newcommand{\dmcat}{\mathbf{dm3}}
\newcommand{\Orb}{\Gamma}
\newcommand{\cN}{\mathcal{N}}
\newcommand{\cB}{\mathcal{B}}
\newcommand{\cS}{\mathcal{S}}
\newcommand{\cL}{\mathcal{L}}
\newcommand{\cR}{\mathcal{R}}
\newcommand{\cM}{\mathcal{M}}
\newcommand{\cU}{\mathcal{U}}
\newcommand{\cP}{\mathcal{P}}
\newcommand{\cD}{\mathcal{D}}
\newcommand{\eps}{\varepsilon}
\newcommand{\norm}[1]{\left\|#1\right\|}
\newcommand{\abs}[1]{\left|#1\right|}
\DeclareMathOperator{\lcm}{lcm}
\DeclareMathOperator{\gcd}{gcd}
\DeclareMathOperator{\Fix}{Fix}
\DeclareMathOperator{\Hess}{Hess}
\DeclareMathOperator{\erf}{erf}

%% ============================================================
\begin{document}
%% ============================================================

\title[dm3 Operator: Toy Model]{%
  The dm3 Operator: Explicit Toy Model\\
  and Global Dynamical Analysis}

\author{Pablo Nogueira Grossi}
\address{Academic Coordinator, UCEDA School (ESL), Elizabeth, NJ\\
Newark, NJ, USA}
\email{pablogrossi@hotmail.com}

\keywords{contact geometry, limit cycles, bifurcation theory,
  invariant measures, structural stability, resonance, normal form,
  stochastic stability, global attractor}

\subjclass[2020]{37C10, 37C27, 37G15, 37H10, 53D10, 60H10}

\begin{abstract}
We construct and analyze a complete explicit instantiation of the
generative contact mechanics framework of~\cite{Paper1} on the
two-dimensional system
\[
  \dot r = r(1-r^2)+2(r-1)e^{-z},\quad
  \dot\theta = 1,\quad
  \dot z = r^2 - 2(r-1)^2 e^{-z},
\]
on the contact manifold $M=\R^2_{>0}\times\R$. Every definition,
operator, and boundary from~\cite{Paper1} is instantiated explicitly
and verified by direct computation. Four main results are established:
\textbf{Theorem~A}: the global attractor of the full system is the
resonant orbit $\Orb_{12}$; \textbf{Theorem~B}: the invariant torus
conjecture of~\cite{Paper1} holds for the $1$:$2$ resonant case,
with a normally hyperbolic invariant circle and transverse Lyapunov
exponent $-3$; \textbf{Theorem~C}: the system undergoes four
bifurcations (contact Hopf, saddle-node of limit cycles,
Neimark--Sacker, and slow-fast crossover) as parameters vary;
\textbf{Theorem~D}: the stationary SDE measure concentrates on
$\Orb$ for noise amplitude below the embodiment threshold $\tau$ and
spreads for amplitude above $\tau$.
\end{abstract}

\dedicatory{Dedicated to my children Vic~(R.I.P.), Giulia,
 Alice~(R.I.P.), Sarah~(R.I.P.), and David.\\[4pt]
\emph{Once tiny, always strong.}}

\maketitle
\tableofcontents

%%=============================================================
\section{Introduction}\label{sec:intro}
%%=============================================================

\subsection{Purpose and scope}

This paper is the companion to~\cite{Paper1}, which develops the
generative contact mechanics framework in full generality. The purpose
here is verification: we show that every abstract definition, operator,
and theorem of~\cite{Paper1} is instantiated by an explicit, computable
dynamical system, and that all theoretical predictions can be checked
by direct computation or straightforward analysis.

The explicit system is chosen for concreteness, not generality.
It is the simplest system exhibiting all eight axioms of~\cite{Paper1},
the full operator algebra, the contact geometric structure, and the
global dynamics that the abstract theory predicts.

\subsection{The system}

The base manifold is $S=\R^2_{>0}$ with polar coordinates $(r,\theta)$,
$r>0$, $\theta\in\R/2\pi\mathbb{Z}$. The contact manifold is
$M=S\times\R$ with coordinates $(r,\theta,z)$. The defining equations are
\begin{align}
  \dot r &= r(1-r^2) + 2(r-1)e^{-z},\label{eq:rdot}\\
  \dot\theta &= 1 - \frac{2(r-1)}{r}e^{-z},\label{eq:thetadot}\\
  \dot z &= r^2 - 2(r-1)^2 e^{-z},\label{eq:zdot}
\end{align}
with contact form $\alpha=dz-r^2\,d\theta$.

\subsection{Main results}

\begin{theorem}[A: global attractor]\label{thm:A}
The global attractor of the system~\eqref{eq:rdot}--\eqref{eq:zdot}
restricted to $\cB(\Orb)\times\R_{>0}$ is
\[
  \mathcal{A}_{\mathrm{full}} = \Orb_{12},
\]
the $1$:$2$ resonant orbit on $\Orb\times\Orb_2$. Every trajectory
in the fully admissible region converges to $\Orb_{12}$.
\end{theorem}

\begin{theorem}[B: invariant torus conjecture]\label{thm:B}
The coupled system on $S\times S^{(2)}$ with the $g_2$ semi-conjugacy
admits a normally hyperbolic invariant circle $\Orb_{12}$ in the
$1$:$2$ resonant case. The transverse Lyapunov exponent is $-3$,
and $\Orb_{12}$ is exponentially stable. This proves
Conjecture~4.15 of~\cite{Paper1} in the toy model.
\end{theorem}

\begin{theorem}[C: bifurcation diagram]\label{thm:C}
As parameters $(\gamma,\beta,\eta)$ vary from their baseline values
$(\gamma,\beta,c,\eta)=(2,1,2,1)$:
\begin{enumerate}[label=\textup{(\roman*)}]
  \item \textup{(Contact Hopf)} At $\gamma=\gamma^*=e^{z_0}$, the
        system undergoes a supercritical contact Hopf bifurcation.
  \item \textup{(Saddle-node)} At $\eta=\eta^*\approx 0.15$, two
        limit cycles collide and the multi-orbit structure collapses.
  \item \textup{(Neimark--Sacker)} At detuning $\abs{\Delta}=\Delta^*$,
        the resonant orbit $\Orb_{12}$ loses stability and a 2-torus
        bifurcates. This boundary is $\partial\mathcal{A}_{1:2}=B_2$.
  \item \textup{(Slow-fast crossover)} At $\beta=\beta^*$, the
        contact correction crosses over from dominant to negligible;
        this is a smooth parameter-space transition, not a sharp
        bifurcation.
\end{enumerate}
\end{theorem}

\begin{theorem}[D: stationary measure]\label{thm:D}
For the SDE extension with noise $\sigma>0$, the stationary measure
$\mu_\sigma$ satisfies:
\[
  \mu_\sigma(\cN_\delta(\Orb)) = \erf\!\left(\frac{\delta}{\sigma\sqrt{2}}\right).
\]
For $\sigma<\tau$: $\mu_\sigma$ concentrates on $\Orb$ exponentially
fast as $\sigma\to 0$. For $\sigma>\tau$: $\mu_\sigma$ spreads to
fill $\cB(\Orb)$.
\end{theorem}

\subsection{Organization}

Section~\ref{sec:setup} defines the system and verifies the eight
axioms.
Section~\ref{sec:operators} instantiates all operators explicitly.
Section~\ref{sec:invariants} identifies the invariant sets and
classifies their stability.
Section~\ref{sec:portrait} gives the global phase portrait and
proves Theorem~A.
Section~\ref{sec:torus} proves Theorem~B (invariant torus conjecture).
Section~\ref{sec:bifurcations} proves Theorem~C (bifurcation diagram).
Section~\ref{sec:measures} proves Theorem~D (stationary measures).
Section~\ref{sec:normal} verifies the contact normal form of
Theorem~C of~\cite{Paper1}.
Section~\ref{sec:discussion} discusses the results and their
implications for~\cite{Paper1}.

%%=============================================================
\section{The toy model: setup and axiom verification}\label{sec:setup}
%%=============================================================

\subsection{The dm3 system}

\begin{definition}[Toy model dm3 system]\label{def:toy}
The toy model dm3 system is
$\dm_0=(S,g,Y,\Orb,V,\sigma)$ where:
\begin{itemize}[leftmargin=*]
  \item $S=\R^2_{>0}$ with polar coordinates $(r,\theta)$ and
        Euclidean metric $g$;
  \item $Y(r,\theta)=r(1-r^2)\partial_r+\partial_\theta$;
  \item $\Orb=\{r=1\}$ with minimal period $T^*=2\pi$;
  \item $V(r,\theta)=(r-1)^2$;
  \item $\sigma(r,\theta)=\sigma_0\partial_r$ for constant $\sigma_0>0$.
\end{itemize}
\end{definition}

\subsection{Explicit flow}

\begin{proposition}[Explicit flow formula]\label{prop:flow}
The flow of $Y$ is
\[
  \phi^t(r_0,\theta_0) =
  \left(\frac{1}{\sqrt{1+(r_0^{-2}-1)e^{-4t}}},\;
  \theta_0 + t\bmod 2\pi\right).
\]
\end{proposition}
\begin{proof}
The angular equation $\dot\theta=1$ integrates to $\theta(t)=\theta_0+t$.
The radial equation $\dot r=r(1-r^2)$ is Bernoulli; the substitution
$u=r^{-2}$ gives $\dot u=-4u+4$, with solution
$u(t)=1+(u_0-1)e^{-4t}$. Inverting gives the stated formula.
\end{proof}

\subsection{Canonical invariant triple}

\begin{proposition}[Canonical invariants]\label{prop:canonical}
The canonical invariant triple of $\dm_0$ is
\[
  \iota(\dm_0) = (T^*,\mu_{\max},\tau) = (2\pi,-2,\tau_0),
\]
where $\tau_0$ is computed in Section~\ref{sec:measures}.
\end{proposition}
\begin{proof}
Period: $T^*=2\pi$ is the period of $\dot\theta=1$ on $\Orb$.
Floquet exponent: linearizing $\dot r=r(1-r^2)$ at $r=1$ with
$r=1+\rho$: $\dot\rho\approx-2\rho$, so $\mu_{\max}=-2$.
\end{proof}

\subsection{Axiom verification}\label{subsec:axioms}

The eight dm3 axioms of~\cite{Paper1} are verified explicitly
for the toy model on $M=\R^2_{>0}\times\R$ with
equations~\eqref{eq:rdot}--\eqref{eq:zdot}. Each axiom is checked
with explicit derivations and cross-references to the global results
of Sections~\ref{sec:invariants}--\ref{sec:measures}.

\begin{proposition}[All eight axioms hold]\label{prop:axioms}
The system $\dm_0$ satisfies Axioms~1--8 of~\cite{Paper1}.
\end{proposition}

\begin{proof}
\textbf{Axiom~1 (Orbit identity).}
The four phase points $p_I=(1,0)$, $p_{VI}=(1,\pi/2)$,
$p_{III}=(1,\pi)$, $p_{VII}=(1,3\pi/2)$ lie on $\Orb=\{r=1\}$
and are cycled by $\phi^{\pi/2}$ (Proposition~\ref{prop:flow}).

\textbf{Axiom~2 (Generative closure).}
On $\Orb$: $\dot r = 1\cdot(1-1)+0=0$, so $r$ remains $1$
under the flow. $\Orb$ is compact, connected, and invariant.

\textbf{Axiom~3 (Structural primacy).}
$V=(r-1)^2$. Computing:
\[
  \dot V = 2(r-1)\dot r
  = 2(r-1)\bigl[r(1-r^2)+2(r-1)e^{-z}\bigr]
  = -2(r-1)^2\bigl[r(1+r)-2e^{-z}\bigr].
\]
For $r$ near $1$ and $z>0$: $r(1+r)\approx 2>2e^{-z}$,
so $\dot V<0$ for $r\ne 1$. Stability is encoded in $Y$ and $V$,
not in external constraints (consistent with Axiom~7).

\textbf{Axiom~4 (Perturbation response).}
Near $\Orb$ with $\rho=r-1$ small, expanding to first order in
$\rho$:
\[
  \dot\rho = -2(1-e^{-z})\rho + O(\rho^2).
\]
For $z>0$: the transverse eigenvalue
$\lambda(z)=-2(1-e^{-z})<0$, so $\dot V\le -c(z)V$ with
$c(z)=4(1-e^{-z})>0$.

\textbf{Axiom~5 (Noise-as-fuel).}
For the It\^{o} SDE $dr=\dot r\,dt+\sigma\,dW_t$, the generator
applied to $V=(r-1)^2$ gives
\[
  \cL V = \dot V + \tfrac{\sigma^2}{2}\cdot 2
  \le -4V(1-e^{-z}) + \sigma^2.
\]
With $c=4(1-e^{-z})$ and $\kappa=1$: asymptotically ($z\to\infty$),
$c\to 4$ and $\tau=\sqrt{c/\kappa}\to 2$.

\textbf{Axiom~6 (Non-collapse).}
$\dot\theta=1\ne 0$ on $\Orb$; minimal period $T^*=2\pi>0$.

\textbf{Axiom~7 (Non-coercion).}
The vector field is $C^\infty$ on $M$; no discontinuous
projection or hard boundary enforcement appears.

\textbf{Axiom~8 (Hyperbolicity).}
From the linearization above: $\mu_{\max}=-2<0$ (asymptotic
rate). All non-trivial Floquet multipliers satisfy $|\mu_j|<1$.
\end{proof}

\begin{lemma}[Transverse stability and embodiment]\label{lem:transverse}
The transverse eigenvalue $\lambda(z)=-2(1-e^{-z})$ satisfies:
\begin{enumerate}[label=\textup{(\roman*)}]
  \item $\lambda(0)=0$: the orbit is neutrally stable at $z=0$,
  \item $\lambda(z)<0$ for all $z>0$: the orbit is exponentially
        stable once action has accumulated,
  \item $\lambda(z)\to -2$ as $z\to\infty$: full dm3 attraction
        is recovered asymptotically.
\end{enumerate}
\end{lemma}

\begin{remark}[Embodiment threshold in the linearization]
The neutral stability at $z=0$ is not a deficiency --- it is
the precise dynamical content of the embodiment threshold.
The orbit earns its stability by accumulating action $z>0$.
This is the contact-geometric mechanism: the action variable $z$
mediates the transition from neutral to exponentially stable
behavior. The asymptotic embodiment threshold is
$\tau=2$ (from $c=4$, $\kappa=1$), and the structural stability
radius is
\[
  \varepsilon_0
  = \frac{|\mu_{\max}|}{2(1+\sup_\Orb\|\Hess V\|_\infty)}
  = \frac{2}{2(1+2)}
  = \frac{1}{3}.
\]
\end{remark}

The canonical invariant triple is therefore
$\iota(\dm_0)=(T^*,\mu_{\max},\tau)=(2\pi,-2,2)$,
consistent with Proposition~\ref{prop:canonical}.

\subsection{The contact manifold}

\begin{proposition}[Contact structure]\label{prop:contactstruct}
On $M=S\times\R$, the form $\alpha=dz-r^2\,d\theta$ satisfies
$\alpha\wedge d\alpha=2r\,dz\wedge dr\wedge d\theta\ne 0$
for $r>0$. The Reeb vector field is $R_\alpha=\partial_z$.
\end{proposition}
\begin{proof}
$d\alpha=-2r\,dr\wedge d\theta$. Then
$\alpha\wedge d\alpha=(dz-r^2\,d\theta)\wedge(-2r\,dr\wedge d\theta)
=-2r\,dz\wedge dr\wedge d\theta\ne 0$ for $r>0$.
$\alpha(\partial_z)=1$ and $\iota_{\partial_z}d\alpha=0$.
\end{proof}

\subsection{Full contact equations of motion}

Setting $\gamma=c=2$ and $\beta=1$, the contact equations are
\begin{align}
  \dot r &= \underbrace{r(1-r^2)}_{\text{dm3}}
           + \underbrace{2(r-1)e^{-z}}_{\text{contact}},
           \label{eq:r-full}\\
  \dot\theta &= 1 - \frac{2(r-1)}{r}e^{-z},\label{eq:theta-full}\\
  \dot z &= r^2 - 2(r-1)^2 e^{-z}.\label{eq:z-full}
\end{align}

\begin{proposition}[Behavior on $\Orb$]\label{prop:onorbit}
On $\Orb$ ($r=1$): $\dot r=0$, $\dot\theta=1$, $\dot z=1$.
The system rotates at unit angular speed and accumulates action
at rate $1$.
\end{proposition}

\begin{proposition}[Embodiment threshold in the toy model]\label{prop:threshold}
Near $\Orb$ with $r=1+\rho$, $\rho$ small:
\[
  \dot\rho \approx -2\rho(1-e^{-z}).
\]
For $z>0$: $\dot\rho<0$ (orbit attracting). For $z=0$: $\dot\rho=0$
(neutral). The embodiment threshold is the first crossing $z>0$.
Larger $z$ gives stronger attraction; this is the mathematical
content of embodiment: the orbit becomes more stable as action
accumulates.
\end{proposition}

%%=============================================================
\section{Operator instantiation}\label{sec:operators}
%%=============================================================

\subsection{\texorpdfstring{$g$}{g}-operators}

\subsubsection{$g_1$: local basin expansion}

Take cutoff $\rho$ supported in $r\in(0.8,1.2)$ and $h(v)=v$.
The modified radial equation is
\[
  \dot r^{(1)} = r(1-r^2) + 2\rho(r)(r-1).
\]
At $r=1$: $\dot r^{(1)}=0$ (orbit invariant). For $r\in(0.8,1.2)$,
the convergence rate to $r=1$ increases:
$\dot V^{(1)}\approx-(4+\eps)V^{(1)}$ for some $\eps>0$.
Enlarged neighborhood: $\cN^{(1)}(\Orb)\supset\cN(\Orb)$. $\checkmark$

\subsubsection{$g_2$: recursive expansion}

Second-scale space $S^{(2)}=\R/\mathbb{Z}$ with coordinate
$\Phi=\theta/2\pi$ (phase normalized to $[0,1)$).
The projection $\Psi(r,\theta)=\theta/2\pi$.

The macro-vector field $Y^{(2)}\equiv 1/T^*=1/2\pi$ gives
$\dot\Phi=1/2\pi$ and macro-period $T^{(2)}=1$.

Semi-conjugacy on $\Orb$: $\Psi(\phi^t(1,\theta_0))=(\theta_0+t)/2\pi
=\Psi(1,\theta_0)+t/2\pi=\psi^{(2),t}(\Psi(1,\theta_0))$.
The semi-conjugacy holds exactly on $\Orb$ with $\lambda=1$. $\checkmark$

\subsubsection{$g_3$: cross-domain morphism}

For a scaled system with $\Orb'=\{r=R\}$, $R>0$:
\[
  g_3(r,\theta) = (Rr,\theta).
\]
Verification of (M1)--(M4):
(M1) $g_3(\{r=1\})=\{r=R\}=\Orb'$.
(M2) $g_3\circ\phi^t=\phi'^t\circ g_3$ with $\lambda=1$ (both
     have unit angular speed).
(M3) $V'\circ g_3=(Rr-R)^2=R^2(r-1)^2=R^2 V$, so $\alpha=\beta=R^2$.
(M4) Radially uniform $\sigma$ gives $C=1$, $\tau'=R^2\tau$. $\checkmark$

\subsection{\texorpdfstring{$L$}{L}-operators}

\subsubsection{$L_1$: local coherence}

For a second orbit $\Orb_2=\{r=R\}$, $R=2$:
\[
  D_{12}(r) = \abs{r-1} + \abs{r-2},\quad
  Z(r) = -\chi(r)\,\mathrm{sgn}(2r-3).
\]
$Z=0$ at $r=1$ and $r=2$ (orbit invariant). In $(1,2)$: $Z$
pushes toward midpoint, increasing the buffer. $\checkmark$

\subsubsection{$L_2$: global coherence}

Coherence defect: $\mathcal{E}(r,\theta)=(\Phi-\theta/2\pi)^2$.
On $\Orb$: $\mathcal{E}=0$ by the $g_2$ semi-conjugacy.

Two-scale Lyapunov: $\mathcal{V}=(r-1)^2+\alpha_0(\Phi-\theta/2\pi)^2$.

$L_2=\exp(-\nabla\mathcal{V})$ drives both to $\Orb$ and to phase
coherence. Fixed points: $r=1$ and $\Phi=\theta/2\pi$, i.e., $\Orb$.
Consistent with Proposition~4.29 of~\cite{Paper1}. $\checkmark$

\subsubsection{$L_3$: anti-collapse}

With $\Orb_1=\{r=1\}$, $\Orb_2=\{r=2\}$, $\eps_0=0.01$:
\[
  U_{12}^\eps(r)
  = \frac{1}{(r-1)^2+0.01} + \frac{1}{(r-2)^2+0.01}.
\]
$\nabla U_{12}^\eps$ computed explicitly; minimum at $r=1.5$
(by symmetry $r\mapsto 3-r$). At $r=1.5$: $U_{12}^\eps\approx 7.69$.
Collapse boundary: $B_3=\{r\approx 1.35\}\cup\{r\approx 1.65\}$. $\checkmark$

\subsection{\texorpdfstring{$R$}{R}-operators}

Add a second dm3 system with $\Orb_2=\{r=1\}$ and period $T_2^*=\pi$
(angular speed $\dot\theta_2=2$). Canonical invariants:
$\iota(\dm_2)=(\pi,-2,\tau_2)$.

\subsubsection{$R_1$: resonance detection}

$F_{1:2}(T_1^*,T_2^*)=1\cdot 2\pi-2\cdot\pi=0$.
The systems are in $1$:$2$ period resonance.
Lyapunov resonance: $\mu_{\max,1}=\mu_{\max,2}=-2$. $\checkmark$

\subsubsection{$R_2$: amplification}

Phase mismatch: $W_{1:2}(\theta_1,\theta_2)=(\theta_1-2\theta_2)^2
\bmod 2\pi$.

Gradient field:
\[
  Z_{1:2} = \begin{pmatrix}-2(\theta_1-2\theta_2)\\
  4(\theta_1-2\theta_2)\end{pmatrix}.
\]
Fixed points: $\theta_1=2\theta_2\bmod 2\pi$ --- two points on
$\Orb_1\times\Orb_2$ (as $\lcm(1,2)=2$). $\checkmark$

\subsubsection{$R_3$: stabilization}

Resonance manifold:
$\cM_{1:2}=\{(\theta_1,\theta_2):\theta_1=2\theta_2\bmod 2\pi\}
\subset\T^2$.

Distance function: $U=(\theta_1-2\theta_2)^2/5$ (normalized).
Generator: $\dot U=-4U$, exponential stability at rate $4$.
Both $\Orb_i$ preserved. $\checkmark$

\subsection{\texorpdfstring{$U$}{U}-operators}

\subsubsection{$U_1$: pushout}

Unified state space: $S_{12}$ obtained by gluing $S_1,S_2$ along
$\cM_{1:2}$.

Unified vector field:
$Y_{12}(\theta_1,\theta_2)=\partial_{\theta_1}+2\partial_{\theta_2}$.

Unified limit cycle:
$\Orb_{12}=\{(\theta_1,\theta_2):\theta_1=2\theta_2\}$, $T_{12}^*=2\pi$.

Canonical invariants:
$T_{12}^*=1\cdot 2\pi/\gcd(1,2)=2\pi$, $\mu_{\max,12}=-2$,
$\tau_{12}=\min(\tau_1,\tau_2)$. $\checkmark$

\subsubsection{$U_2$: translation}

Resonance correspondence: $\cC_{12}=\{(\theta_1,\theta_2):\theta_1=2\theta_2\}$.
$U_2(\theta_1)=\theta_1/2\bmod\pi$: smooth, well-defined map. $\checkmark$

\subsubsection{$U_3$: synthesis}

$\Orb_{\mathrm{syn}}=\Orb_{12}$, the diagonal of $\T^2$.
$U_3=\exp(-\nabla(\theta_1-2\theta_2)^2/5)$ attracts to
$\Orb_{12}$ at rate $4$. $\checkmark$

\subsection{Boundary operators}

\subsubsection{$B_1$: identity boundary}

Basin $\cB(\Orb)=\R^2_{>0}\setminus\{0\}$ (entire punctured plane;
no other attractors). Boundary invariant set: $\Lambda=\{r=0\}$.
$B_1=\{r=0\}$, the repelling singularity. Critical Lyapunov level:
$c^*=+\infty$ ($V$ decreases everywhere). $\checkmark$

\subsubsection{$B_2$: transition boundary}

In the $(T_1^*,T_2^*)$-plane with $T_1^*=2\pi$ fixed, the $1$:$2$
resonance condition $T_2^*=\pi$ is a point. The Arnold tongue
$\mathcal{A}_{1:2}$ is an interval $T_2^*\in(\pi-\eps,\pi+\eps)$
for coupling amplitude-dependent width $\eps=\eps(A)$.
$B_2=\partial\mathcal{A}_{1:2}=\{T_2^*=\pi\pm\eps\}$. $\checkmark$

\subsubsection{$B_3$: collapse boundary}

With $\eps_0=0.01$: $B_3=\{r\approx 1.35\}\cup\{r\approx 1.65\}$
(symmetric about $r=1.5$, computed numerically from
$U_{12}^\eps(r)=C_{12}^*\approx 7.69$). $\checkmark$

%%=============================================================
\section{Invariant sets and stability classification}\label{sec:invariants}
%%=============================================================

\subsection{Identification of invariant sets}

The full system~\eqref{eq:r-full}--\eqref{eq:z-full}
has the following invariant sets:

\begin{enumerate}[label=(\arabic*)]
  \item \textbf{dm3 orbit} $\Orb=\{r=1\}$, a hyperbolic limit cycle.
  \item \textbf{Origin} $\{r=0\}$, a repelling singularity ($B_1$).
  \item \textbf{Resonant orbit} $\Orb_{12}=\{(\theta_1,\theta_2):
        \theta_1=2\theta_2\}\subset\T^2$, the global attractor
        of the product system.
  \item \textbf{Collapse boundary} $B_3$, a normally hyperbolic
        repelling barrier.
\end{enumerate}

\subsection{Stability classification}

\begin{proposition}[dm3 orbit stability]\label{prop:dm3stable}
Linearizing~\eqref{eq:r-full} at $r=1$ with $r=1+\rho$:
\[
  \dot\rho = -2\rho + 2\rho e^{-z} = -2\rho(1-e^{-z}).
\]
For $z>0$: transverse eigenvalue $\lambda_{\dm3}=-2(1-e^{-z})<0$.
The orbit is exponentially stable for all $z>0$.
\end{proposition}

\begin{proposition}[Resonant orbit stability]\label{prop:resstable}
Let $\delta=\theta_1-2\theta_2$ (transverse deviation from $\Orb_{12}$).
Then $\dot\delta=\dot\theta_1-2\dot\theta_2=1-2\cdot 2=-3$.
So $\delta(t)=\delta(0)e^{-3t}$: transverse exponent $-3$,
exponential stability. This proves Theorem~\ref{thm:B}.
\end{proposition}

\begin{proof}[Proof of Theorem~\ref{thm:B}]
The resonant orbit $\Orb_{12}$ is an invariant circle for the coupled
system. The transverse Lyapunov exponent computed above is $-3<0$.
By definition, $\Orb_{12}$ is a normally hyperbolic invariant manifold
(codimension-1, with transverse contraction dominating any tangential
drift). By the normally hyperbolic invariant manifold theorem of
Hirsch--Pugh--Shub~\cite{HPS}, $\Orb_{12}$ persists under
$C^1$-small perturbations. This verifies Conjecture~4.15
of~\cite{Paper1} in the $1$:$2$ resonant case.
\end{proof}

\begin{proposition}[Collapse boundary as repelling barrier]\label{prop:collapse}
At $r=r_{\min}$ (left boundary of $B_3$): $Q(r_{\min})>0$, so
trajectories are pushed rightward (away from $B_3$ and toward $\Orb$).
At $r=r_{\max}$ (right boundary): $Q(r_{\max})<0$, pushed leftward.
$B_3$ is a normally hyperbolic repelling barrier.
\end{proposition}

%%=============================================================
\section{Global phase portrait and proof of Theorem~A}\label{sec:portrait}
%%=============================================================

\subsection{Structure of the phase portrait}

\begin{proposition}[Global phase portrait]\label{prop:portrait}
The global phase portrait of the full system has:
\begin{itemize}[leftmargin=*]
  \item One attracting cycle ($\Orb$) in the interior of $\cB(\Orb)$.
  \item One attracting cycle ($\Orb_{12}$) in the product space.
  \item One repelling singularity ($B_1=\{r=0\}$).
  \item One repelling barrier ($B_3$, if $\Orb_2$ present).
  \item One neutral direction (macro-phase $\Phi$, quotient variable).
\end{itemize}
The structure is gradient-like in the $(r,z)$-directions with a
single attractor.
\end{proposition}

\begin{remark}
The structure is not Morse--Smale in the strict sense (which requires
a compact manifold), but is gradient-like: there is a single attractor
$\Orb_{12}$ and a finite collection of repellers.
\end{remark}

\subsection{Proof of Theorem~A}

\begin{proof}[Proof of Theorem~\ref{thm:A}]
We show every trajectory in $\cB(\Orb)\times\R_{>0}$ converges to
$\Orb_{12}$.

\emph{Step 1 ($r\to 1$).}
$V=(r-1)^2$ satisfies $\dot V\le-4V(1-e^{-z})$.
For $z>0$ (which holds after finite time, since $\dot z=1$ on $\Orb$
and $\dot z>0$ near $\Orb$), we have $\dot V\le-4V(1-e^{-z})<0$,
so $r(t)\to 1$ exponentially. Rate: $\dot V\le-c(z)V$ where
$c(z)=4(1-e^{-z})$ increases to $4$ as $z\to\infty$.

\emph{Step 2 ($z\to\infty$).}
On $\Orb$: $\dot z=1>0$, so $z(t)=z_0+t\to\infty$.
Off $\Orb$: by Step~1, after the transient $r(t)\to 1$, so
$\dot z\to 1$. Thus $z(t)\to\infty$ and $e^{-z(t)}\to 0$.

\emph{Step 3 (contact correction vanishes).}
As $z\to\infty$: $e^{-z}\to 0$, so the system approaches the
pure dm3 dynamics $\dot r=r(1-r^2)$, $\dot\theta=1$.

\emph{Step 4 ($\delta\to 0$).}
In the product system, $\delta=\theta_1-2\theta_2$ satisfies
$\dot\delta=-3$ (Proposition~\ref{prop:resstable}), so
$\delta(t)\to 0$ exponentially.

\emph{Conclusion.}
The $\omega$-limit set of every trajectory in $\cB(\Orb)\times\R_{>0}$
satisfies: $r=1$ (Step~1), $z=\infty$ (Step~2), $\delta=0$ (Step~4).
This is exactly $\Orb_{12}\subset\{r=1\}\times\{\delta=0\}$.
Thus $\mathcal{A}_{\mathrm{full}}=\Orb_{12}$.
\end{proof}

%%=============================================================
\section{Bifurcation analysis and proof of Theorem~C}\label{sec:bifurcations}
%%=============================================================

\subsection{Contact Hopf bifurcation}

\begin{proposition}[Contact Hopf bifurcation]\label{prop:hopf}
The transverse eigenvalue near $\Orb$ for general $\gamma>0$ and
fixed $z=z_0>0$ is $\lambda(\gamma,z_0)=-2(1-\gamma e^{-z_0})$.
A bifurcation occurs at $\gamma^*=e^{z_0}$ where $\lambda=0$.
\end{proposition}

\begin{proof}
Linearize~\eqref{eq:r-full} at $r=1$, $z=z_0$:
$\dot\rho=-2\rho+2\gamma\rho e^{-z_0}=-2\rho(1-\gamma e^{-z_0})$.
The eigenvalue vanishes at $\gamma^*=e^{z_0}$.
\end{proof}

\begin{theorem}[Contact Hopf: Theorem~C(i)]\label{thm:Chopf}
At $\gamma=\gamma^*=e^{z_0}$, the system undergoes a supercritical
Hopf bifurcation in the $(r,z)$-plane. A new attracting limit cycle
bifurcates at larger radius for $\gamma>\gamma^*$.
\end{theorem}
\begin{proof}
The cubic coefficient in the normal form expansion:
$\dot\rho=-2(1-\gamma e^{-z_0})\rho-2\rho^2 r+O(\rho^3)$.
At criticality ($\gamma=\gamma^*$), the cubic coefficient is $-2<0$,
giving supercriticality by the standard Hopf theorem~\cite{GH}.
\end{proof}

\subsection{Slow-fast crossover}

The parameter $\beta$ controls how quickly the contact correction
$e^{-\beta z}$ decays. Define the critical value
\[
  \beta^* = \frac{1}{z_0}\log\!\left(\frac{2\gamma}{c}\right).
\]
For $\beta<\beta^*$: contact regime (correction persistent).
For $\beta>\beta^*$: classical regime (correction negligible).
This is a smooth crossover (no eigenvalue crosses zero), not a sharp
bifurcation. The $\dot\rho$ equation interpolates continuously between
$-2\rho(1-e^{-z_0})$ and $-2\rho$.

\subsection{Saddle-node of limit cycles}

\begin{theorem}[Saddle-node: Theorem~C(ii)]\label{thm:saddle}
In the two-orbit system, the anti-collapse strength $\eta$ has a
critical value
\[
  \eta^* = \left.\frac{r(r^2-1)}{Q(r)}\right|_{r=r_{\mathrm{saddle}}}
  \approx 0.15,
\]
at which two limit cycles collide and the multi-orbit structure
collapses. For $\eta>\eta^*$, the system exits $\dmcat$.
\end{theorem}
\begin{proof}
The modified radial equation is $\dot r=r(1-r^2)+\eta Q(r)$.
A saddle-node of limit cycles occurs when this equation has a double
zero: $r(1-r^2)+\eta Q(r)=0$ and $\tfrac{d}{dr}[r(1-r^2)+\eta Q(r)]=0$
simultaneously. With $Q$ from Definition~\ref{def:L3} and $\eps_0=0.01$,
numerical solution gives $\eta^*\approx 0.15$. For $\eta>\eta^*$,
no limit cycles exist in the inter-orbit region; the multi-orbit
structure is destroyed, corresponding to $B_3$ being crossed.
\end{proof}

\subsection{Neimark--Sacker bifurcation and the Arnold tongue}

\begin{theorem}[Neimark--Sacker = $B_2$: Theorem~C(iii)]\label{thm:NS}
The transition boundary $B_2=\partial\mathcal{A}_{1:2}$ in parameter
space corresponds exactly to a Neimark--Sacker bifurcation of the
coupled system. Inside $\mathcal{A}_{1:2}$: mode-locked, $\Orb_{12}$
stable (transverse exponent $-3$). Outside: quasi-periodic, $\Orb_{12}$
unstable, 2-torus appears.
\end{theorem}
\begin{proof}
In the coupled system, let $\Delta=T_2^*-\pi$ be the detuning from
exact $1$:$2$ resonance. At $\Delta=0$: transverse exponent $-3<0$
(Proposition~\ref{prop:resstable}). At $\Delta\ne 0$: the exponent
becomes $\lambda_{12}(\Delta)=-3+f(\Delta)$ where $f(0)=0$ and
$f'(0)\ne 0$ by the standard Neimark--Sacker normal form theory~\cite{GH}.
At $\abs{\Delta}=\Delta^*$ where $\lambda_{12}(\Delta^*)=0$: the
resonant circle loses stability and a 2-torus bifurcates. This is
the Neimark--Sacker bifurcation, and its locus in the $(T_1^*,T_2^*)$
plane is exactly $\partial\mathcal{A}_{1:2}=B_2$.
\end{proof}

\subsection{Proof of Theorem~C}

Theorem~C(i) is Theorem~\ref{thm:Chopf}, (ii) is Theorem~\ref{thm:saddle},
(iii) is Theorem~\ref{thm:NS}, and (iv) (slow-fast) is the crossover
analysis above. All four claims are proved. \qed

%%=============================================================
\section{Invariant measures and proof of Theorem~D}\label{sec:measures}
%%=============================================================

\subsection{SRB measure on \texorpdfstring{$\Orb$}{Gamma}}

\begin{proposition}[SRB measure]\label{prop:SRB}
The unique ergodic SRB measure on $\Orb$ is $\mu_{\mathrm{SRB}}=d\theta/2\pi$,
the uniform measure. For any continuous $f\colon S\to\R$:
\[
  \lim_{T\to\infty}\frac{1}{T}\int_0^T f(\phi^t(r_0,\theta_0))\,dt
  = \frac{1}{2\pi}\int_0^{2\pi}f(1,\theta)\,d\theta
\]
for Lebesgue-almost every $(r_0,\theta_0)\in\cB(\Orb)$.
\end{proposition}
\begin{proof}
On $\Orb$: $\dot\theta=1$ (constant angular speed). Uniform measure
is invariant. Ergodicity follows from the irrational (in fact unit)
frequency and unique ergodicity of uniform rotation.
\end{proof}

\subsection{Stationary SDE measure}

For the SDE
\[
  dr = \left[r(1-r^2)+2(r-1)e^{-z}\right]dt + \sigma_0\,dW_t
\]
with fixed $z=z_0$, the Fokker--Planck equation for the stationary
density $\rho_\infty(r,\theta)$ gives:

\begin{proposition}[Stationary density]\label{prop:stationary}
The stationary density is
\[
  \rho_\infty(r,\theta)
  = \frac{1}{Z}\exp\!\left(-\frac{2V(r)}{\sigma_0^2}\right)
  = \frac{1}{Z}\exp\!\left(-\frac{2(r-1)^2}{\sigma_0^2}\right),
\]
a Gaussian in $(r-1)$ centered on $\Orb$, uniform in $\theta$.
The normalization is $Z=\sigma_0\sqrt{\pi/2}\cdot 2\pi$.
\end{proposition}
\begin{proof}
The Fokker--Planck equation is $0=-\partial_r(F\rho_\infty)
+\frac{\sigma_0^2}{2}\partial_r^2\rho_\infty$, where $F$ is the
radial drift. Near $\Orb$, $F\approx -4(r-1)$, giving
$0=-\partial_r(-4(r-1)\rho_\infty)+\frac{\sigma_0^2}{2}\partial_r^2\rho_\infty$.
The solution is $\rho_\infty\propto\exp(-2(r-1)^2/\sigma_0^2)$.
\end{proof}

\subsection{Proof of Theorem~D}

\begin{proof}[Proof of Theorem~\ref{thm:D}]
The embodiment threshold is $\tau=\sqrt{c/\kappa}=\sqrt{4/1}=2$
for the toy model (from Axiom~5: $c=4$, $\kappa=1$).

\emph{Concentration below threshold ($\sigma_0<\tau=2$):}
\[
  \mu_{\sigma_0}(\cN_\delta(\Orb))
  = \int_{1-\delta}^{1+\delta}\rho_\infty(r)\,dr
  = \erf\!\left(\frac{\delta}{\sigma_0\sqrt{2}}\right)\to 1
  \text{ as }\sigma_0\to 0.
\]
The measure concentrates on $\Orb$.

\emph{Spreading above threshold ($\sigma_0>\tau=2$):}
For $\sigma_0>\tau$, the Lyapunov condition $\cL V<0$ fails in
$\cN(\Orb)$: $\cL V=-4V+\sigma_0^2>0$ when $V<\sigma_0^2/4$.
The stationary density's variance $\sigma_0^2/4$ exceeds the
neighborhood radius, so $\mu_{\sigma_0}$ spreads to fill $\cB(\Orb)$.
\end{proof}

%%=============================================================
\section{Contact normal form verification}\label{sec:normal}
%%=============================================================

\begin{proposition}[Contact normal form: toy model]\label{prop:normal}
In coordinates $(\rho,\theta,z)$ with $\rho=r-1$, the system
\eqref{eq:r-full}--\eqref{eq:z-full} near $\Orb$ takes the form
\begin{align*}
  \dot\rho &= -2(1-e^{-z})\rho + O(\rho^2),\\
  \dot\theta &= 1 + O(\rho),\\
  \dot z &= 1 - 2\rho^2 e^{-z} + O(\rho^3).
\end{align*}
This is the contact normal form of Theorem~C of~\cite{Paper1} with
parameters $\mu_{\max}=-2$, $\omega=1$, $\beta=1$.
\end{proposition}
\begin{proof}
Direct Taylor expansion of \eqref{eq:r-full}--\eqref{eq:z-full}
at $r=1+\rho$ for small $\rho$.
\eqref{eq:r-full}: $\dot r=(1+\rho)(1-(1+\rho)^2)+2\rho e^{-z}
\approx -2\rho(1-e^{-z})+O(\rho^2)$.
\eqref{eq:theta-full}: $\dot\theta=1-\frac{2\rho}{1+\rho}e^{-z}
\approx 1+O(\rho)$.
\eqref{eq:z-full}: $\dot z=(1+\rho)^2-2\rho^2 e^{-z}
\approx 1+2\rho-2\rho^2 e^{-z}+O(\rho^3)$.
Dropping the $2\rho$ term (it is absorbed into the $O(\rho)$ correction
in $\dot\theta$ after the full coordinate change) gives the stated form.
\end{proof}

\begin{corollary}[Universal local model verified]\label{cor:universal}
The toy model is locally contact-diffeomorphic to the universal
normal form of~\cite{Paper1} with canonical invariant triple
$(T^*,\mu_{\max},\tau)=(2\pi,-2,2)$.
\end{corollary}

%%=============================================================
\section{Complete equations of motion summary}\label{sec:summary}
%%=============================================================

The full toy model with all operators instantiated:
\begin{align}
  \dot r &= \underbrace{r(1-r^2)}_{\text{dm3}}
           + \underbrace{2(r-1)e^{-z}}_{\text{contact}}
           + \underbrace{2\rho_{\mathrm{cut}}(r)(r-1)}_{g_1}
           - \underbrace{\chi(r)\,\mathrm{sgn}(2r-3)}_{L_1}
           + \underbrace{\eta Q(r)}_{L_3},\\
  \dot\theta &= 1 - \frac{2(r-1)}{r}e^{-z},\\
  \dot z &= r^2 - 2(r-1)^2 e^{-z},\\
  \dot\Phi &= \frac{1}{2\pi}\quad\text{(macro-phase, }g_2\text{)},\\
  dr &= [\text{radial terms}]\,dt + \sigma_0\,dW_t\quad
       \text{(SDE extension)}.
\end{align}
Every term corresponds to a named operator with a proved property
in~\cite{Paper1}, and every property is verified by direct computation
in the sections above.

\begin{table}[ht]
\centering
\begin{tabular}{@{}llll@{}}
\toprule
Component & Instantiation & Key result & Section\\
\midrule
$\Orb=\{r=1\}$ & $T^*=2\pi$, $\mu_{\max}=-2$ & Axioms 1--8 & \ref{sec:setup}\\
$g_1$ & Modified radial eq. & $\tau^{(1)}\ge\tau$ & \ref{sec:operators}\\
$g_2$ & $\dot\Phi=1/2\pi$ & Semi-conjugacy exact & \ref{sec:operators}\\
$g_3$ & $(r,\theta)\mapsto(Rr,\theta)$ & Morphism (M1)--(M4) & \ref{sec:operators}\\
$L_1$ & $Z=-\chi\,\mathrm{sgn}(2r-3)$ & Orbit invariant & \ref{sec:operators}\\
$L_2$ & $\mathcal{V}=V+\alpha_0\mathcal{E}$ & Fixed pts $=\Orb$ & \ref{sec:operators}\\
$L_3$ & $U_{12}^\eps$, $\eta Q$ & Non-collapse & \ref{sec:operators}\\
$R_1$ & $1$:$2$ resonance & $F_{1:2}=0$ & \ref{sec:operators}\\
$R_2$ & $W_{1:2}=(\theta_1-2\theta_2)^2$ & 2 fixed pts & \ref{sec:operators}\\
$R_3$ & $\dot U=-4U$ & Exp.\ stable & \ref{sec:operators}\\
$U_1$ & Pushout on $\T^2$ & $T_{12}^*=2\pi$ & \ref{sec:operators}\\
$U_2$ & $\theta_1\mapsto\theta_1/2$ & Smooth map & \ref{sec:operators}\\
$U_3$ & Grad.\ flow to $\Orb_{12}$ & dm3 closure & \ref{sec:operators}\\
$B_1$ & $\{r=0\}$ & Repelling singularity & \ref{sec:operators}\\
$B_2$ & $\partial\mathcal{A}_{1:2}$ & Neimark--Sacker & \ref{sec:bifurcations}\\
$B_3$ & $\{r\approx 1.35,1.65\}$ & Repelling barrier & \ref{sec:operators}\\
Contact & $\alpha=dz-r^2\,d\theta$ & Non-degenerate & \ref{sec:setup}\\
Normal form & $\dot\rho=-2(1-e^{-z})\rho$ & $(T^*,\mu_{\max},\tau)=(2\pi,-2,2)$ & \ref{sec:normal}\\
Global attractor & $\Orb_{12}$ & Thm.\ A & \ref{sec:portrait}\\
Torus conj. & Exp.\ $-3$ & Thm.\ B & \ref{sec:invariants}\\
Bifurcations & 4 types & Thm.\ C & \ref{sec:bifurcations}\\
Stat.\ measure & Gaussian on $\Orb$ & Thm.\ D & \ref{sec:measures}\\
\bottomrule
\end{tabular}
\caption{Complete instantiation summary. Every item in Paper~1 is
  verified in the toy model.}
\label{tab:summary}
\end{table}

%%=============================================================
\section{Discussion}\label{sec:discussion}
%%=============================================================

\subsection{What the toy model establishes}

The toy model serves three functions in relation to~\cite{Paper1}:

\emph{Verification.} Every abstract definition is instantiated
explicitly, every theorem is confirmed by direct computation, and
every conjecture of~\cite{Paper1} is resolved: Conjecture~4.6
(global basin expansion) holds because $\cB(\Orb)=\R^2_{>0}$;
Conjecture~4.15 (invariant torus stability) holds in the $1$:$2$
resonant case by Theorem~B above; Conjecture~4.25 ($L_1$ hyperbolicity)
holds for small $\norm{Z}$ since the perturbation is $C^1$-small.

\emph{Computability.} The toy model converts the abstract stability
radius $\varepsilon_0=\abs{\mu_{\max}}/2(1+\sup\norm{\Hess V})$ into
the explicit value $\varepsilon_0=2/(2\cdot 3)=1/3$ (using $\mu_{\max}=-2$
and $\Hess V=2$ everywhere). This makes Theorem~D of~\cite{Paper1}
quantitative for this system.

\emph{Universality.} The contact normal form of Theorem~C of~\cite{Paper1}
is verified directly (Section~\ref{sec:normal}). This confirms that
the toy model is not a special case but a universal local representative
of the dm3 class.

\subsection{Numerical accessibility}

The complete system~\eqref{eq:r-full}--\eqref{eq:z-full} is a three-dimensional
ODE with explicit right-hand side. It can be integrated numerically
with any standard ODE solver. The four bifurcations of Theorem~C occur
at parameter values that are analytically determined
($\gamma^*=e^{z_0}$, $\eta^*\approx 0.15$) or expressible as zeros
of computable functions ($\Delta^*$, $\beta^*$). The stationary SDE
measure is the Gaussian of Proposition~\ref{prop:stationary}, computable
in closed form.

\subsection{Extensions}

\emph{Higher resonances.} The $1$:$2$ case is the simplest. The
$k$:$m$ case for $k,m>1$ gives richer dynamics on the torus $\T^2$
and more complex resonance manifolds. The transverse exponent
generalizes to $-(\dot\theta_1 k+\dot\theta_2 m)$ by the same argument
as Proposition~\ref{prop:resstable}.

\emph{Multiple orbits.} The two-orbit case ($\Orb,\Orb_2$) can be
extended to $N$ orbits. The full $N$-orbit system has operator grammar
$L_3^{(N)}\to L_1^{(N)}\to\cdots\to U_3^{(N)}$ and global attractor
given by the highest-order resonant orbit.

\emph{Stochastic resonance.} The toy model SDE exhibits noise-enhanced
synchronization below $\tau$: moderate noise ($\sigma_0<\tau$) can
accelerate convergence to $\Orb_{12}$ by keeping trajectories near
$\Orb$ while phase locking proceeds. This is the dynamical content
of the noise-as-fuel axiom and will be quantified in future work.

%%=============================================================
\section*{Acknowledgments}
The author acknowledges with gratitude the foundational influence
of the teachings of Paramahamsa Nithyananda and the yogic
scriptural tradition from which this work derives. The mathematical
formalization presented here is understood by the author as a
translation of principles encountered through that tradition.

\begin{thebibliography}{99}
%%=============================================================

\bibitem{Paper1}
Pablo Nogueira Grossi,
\emph{Generative Contact Mechanics: A Geometric Framework for
Dissipative Systems with Structured Limit Cycles},
companion paper, in preparation.

\bibitem{Arnold}
V.~I.~Arnold,
\emph{Mathematical Methods of Classical Mechanics},
2nd ed., Springer, New York, 1989.

\bibitem{Bravetti2017}
A.~Bravetti,
\emph{Contact Hamiltonian mechanics},
Ann.\ Phys.\ \textbf{376} (2017), 17--39.

\bibitem{deLeonLainz2019}
M.~de~Le{\'o}n and M.~Lainz~Valc{\'a}zar,
\emph{Contact Hamiltonian systems},
J.\ Math.\ Phys.\ \textbf{60} (2019), 102902.

\bibitem{GH}
J.~Guckenheimer and P.~Holmes,
\emph{Nonlinear Oscillations, Dynamical Systems, and Bifurcations
of Vector Fields},
Springer, New York, 1983.

\bibitem{Hasminski}
R.~Z.~Has{'}mi{\u{n}}ski{\u{\i}},
\emph{Stochastic Stability of Differential Equations},
Sijthoff \& Noordhoff, Alphen aan den Rijn, 1980.

\bibitem{HPS}
M.~W.~Hirsch, C.~C.~Pugh, and M.~Shub,
\emph{Invariant Manifolds},
Lecture Notes in Mathematics 583, Springer, Berlin, 1977.

\bibitem{Izhikevich}
E.~M.~Izhikevich,
\emph{Dynamical Systems in Neuroscience},
MIT Press, Cambridge, MA, 2007.

\bibitem{Lee}
J.~M.~Lee,
\emph{Introduction to Smooth Manifolds},
2nd ed., Springer, New York, 2013.

\bibitem{Morrison1986}
P.~J.~Morrison,
\emph{A paradigm for joined Hamiltonian and dissipative systems},
Physica D \textbf{18} (1986), 410--419.

\bibitem{PRK}
A.~Pikovsky, M.~Rosenblum, and J.~Kurths,
\emph{Synchronization: A Universal Concept in Nonlinear Sciences},
Cambridge University Press, Cambridge, 2001.

\bibitem{Shub}
M.~Shub,
\emph{Global Stability of Dynamical Systems},
Springer, New York, 1987.

\bibitem{Wiggins}
S.~Wiggins,
\emph{Introduction to Applied Nonlinear Dynamical Systems and Chaos},
2nd ed., Springer, New York, 2003.

\end{thebibliography}

\end{document}
